\section{Conclusion}
\label{sec:conclusion}

Published estimates of accelerator depreciation disagree because they read different prices
through different assumptions. In a panel of posted rents for seven NVIDIA chips, realized
rents fall slowly; the shape of the age profile flips between spot and list prices; and with valid
few-cluster inference no identification strategy pins down the rate or the shape. Rents stay far
above operating cost for the whole observed life, so the economic life that matters for book
depreciation is decided by obsolescence in power-constrained data centers, not by rents reaching
operating cost.

\ajcheck{For accounting, the implication is narrower than either side of the public debate. Five-
or six-year straight-line lives are not contradicted by posted rents, whose value curves decline
more slowly than straight-line. They look long only if one believes that power-constrained
operators retire chips when a better chip arrives, which is what the retirement model and the
largest operators' own behavior suggest; even then the model's median retirement age of
\nPThreeMedian{} years has a P10--P90 range of \nPThreePTen--\nPThreePNinety{} years that contains
those lives. The firm-level effect of that choice is reported in
Section~\ref{sec:accounting}.}

The panel, code and identification range are public, so each of these choices can be revisited as
longer price histories for the Hopper and Blackwell generations accumulate.
